%% 中文对照稿（内部阅读用）——投稿版为英文 main.tex；本稿与其句对句对应。
%% Build: latexmk -xelatex main-zh.tex
\documentclass[11pt]{ctexart}

\usepackage{amsmath}
\usepackage{url}
\usepackage{graphicx}
\usepackage[margin=2.6cm]{geometry}
\usepackage{listings}
\usepackage{tikz}
\usetikzlibrary{patterns,arrows.meta,positioning,decorations.pathreplacing,calc}
\setmonofont{Menlo}

%% ---- Lean listings：与英文版相同的 unicode literate 设置 ----
\lstdefinelanguage{lean}{
  keywords={theorem,def,structure,where,fun,by,exact,intro,obtain,refine,show,
            rcases,rw,omega,ring,calc,have,open,import,namespace,end,variable,abbrev},
  sensitive=true,
  comment=[l]{--},
  morecomment=[s]{/-}{-/},
  literate={≤}{{$\leq$}}1 {≥}{{$\geq$}}1 {∀}{{$\forall$}}1 {∃}{{$\exists$}}1
           {∈}{{$\in$}}1 {×}{{$\times$}}1 {⇒}{{$\Rightarrow$}}1 {→}{{$\to$}}1
           {ℕ}{{$\mathbb{N}$}}1 {ℂ}{{$\mathbb{C}$}}1 {≃}{{$\simeq$}}1
           {∧}{{$\wedge$}}1 {¬}{{$\neg$}}1 {∩}{{$\cap$}}1 {∪}{{$\cup$}}1
           {⟨}{{$\langle$}}1 {⟩}{{$\rangle$}}1 {≠}{{$\neq$}}1 {√}{{$\sqrt{}$}}1
           {Θ}{{$\Theta$}}1 {Ω}{{$\Omega$}}1 {μ}{{$\mu$}}1 {θ}{{$\theta$}}1 {ι}{{$\iota$}}1 {σ}{{$\sigma$}}1
           {□}{{$\square$}}1 {·}{{$\cdot$}}1 {↑}{{$\uparrow$}}1 {∉}{{$\notin$}}1,
}
\lstset{language=lean,basicstyle=\small\ttfamily,columns=fullflexible,
        keywordstyle=\bfseries,commentstyle=\itshape\color{gray},
        breaklines=true,frame=single,framerule=0.2pt,xleftmargin=1em}

\newcommand{\cc}{\mathrm{cc}}
\newcommand{\pw}{\mathrm{pw}}
\newcommand{\tw}{\mathrm{tw}}

\title{你的收缩顺序离最优还有多远？\\
       \large sweep 型张量网络收缩的机器核验绝对下界\\[2pt]
       \normalsize （中文对照稿；投稿版为英文 \texttt{main.tex}，本稿句对句对应）}
\author{匿名作者}
\date{}

\begin{document}
\maketitle

\begin{abstract}
张量网络收缩的代价由收缩顺序决定，而寻找最优顺序（TNCO）是 NP-hard 的，实践因此依赖启发
式搜索。搜索结果只能被\emph{相对地}评判——顺序 $A$ 胜过顺序 $B$——因为不存在可信的
\emph{绝对}标尺去界定它们与真实最优的距离。
我们补上这把标尺：对收缩代价给出机器核验的下界，在 \emph{sweep 型（stem）}序类上无条件
成立，在 Lean~4 + Mathlib 中形式化（约 2{,}800 行，无 \texttt{sorry}，\textbf{零自定义公
理}）；推广到任意收缩树依赖唯一一条显式陈述的已发表假设（晶格上 $\pw=\Theta(\tw)$），是
prose 中的文献推论，本身未被形式化。主要的认证部件是\emph{宽度}下界：bramble 是
pathwidth 的组合对偶证书，经
一个初等的区间 Helly 论证完全自证——据我们所知，这是 Lean 中首个 pathwidth 下界理论。第二
张证书刻意保持初等，面向\emph{给定}参数族回应秩侧漏洞：\emph{Contraction Rigidity} 定理把
一个满秩门配置升级为
非零的割面行列式多项式——证书就是这一条非零性，一个单点求值论证，不含概率——且"能力"前件经
Kronecker 路由约化为组合的跨割面匹配，因此对这样一个族，在该多项式零集之外的任何门参数处，没
有低秩（decision diagram 式）因子分解能穿过割面。这张秩侧证书是框架层面的、居于次要地位：
它作为伴随定理与代价下界并列陈述，位于主定理合取之外、从未复合；它在校准割面上的实例化仅认
证割面尺寸层面的\emph{能力}（见证在门参数上为常量、在任何平衡割面上逐字可得）；且它不覆盖
离散门集上的具体电路。在校准算例
——53 线链的
$53\times53$ 深窗口时空 grid，它与
Sycamore-53 的唯一联系是共享的\emph{割面尺度} $53$——上，证书在\emph{宽度上零间隙}：每个
stem 序宽度 $\ge 54$、在模型的 bag 计数代价语义下总
代价 $\ge 2^{54}$，而我们的显式 sweep 达到宽度 $54$——一个携带 stem 类内最优性机器核验证
明的展示解。
\end{abstract}

%% ===================================================================
\section{引言}\label{sec:intro}

张量网络收缩是量子电路经典模拟的主力工具，而它的代价几乎完全由一个离散选择决定：\emph{收缩
顺序}。选好并不容易——在一般网络上寻找最优顺序是 NP-hard 的~\cite{markov2008simulating}
——于是十年的实践孕育出丰富的启发式生态：hyper-optimized 收缩树~\cite{gray2021hyper}、
big-batch 与切片方法~\cite{pan2022solving}、面向 supremacy 电路的 stem 型顺
序~\cite{huang2020classical}、multi-stem 改进~\cite{swtnc2025}、以及复用感知的变
体~\cite{kalachev2021recursive}。这些工具令人印象深刻，但它们共享同一个方法论盲区：
\textbf{它们支持的所有质量评判都是相对的}。一个新启发式的验证方式，是在基准集上击败前一
个。而无论哪个顺序距离真实最优还有多远，都无从得知——因为没有任何可信的东西从下方兜住最优
值。"state of the art" 的字面含义只是"我们碰巧找到的顺序里最好的那个"（图~\ref{fig:yardstick}）。

\begin{figure}[t]
\centering
\begin{tikzpicture}[scale=0.9, font=\small]
  \draw[-{Stealth}] (0,0) -- (0,4.6) node[left=2pt] {代价};
  \fill[pattern=north east lines, pattern color=red!50] (0.15,0.2) rectangle (6.2,1.0);
  \draw[red!70!black, very thick] (0.15,1.0) -- (6.2,1.0)
    node[pos=0.4, above] {\textbf{认证下界} $L$（机器核验的定理）};
  \node[red!70!black] at (3.2,0.55) {此区域不存在任何收缩顺序};
  \fill[blue!70!black] (2.0,3.9) circle (2pt) node[right=3pt] {顺序 $A$};
  \fill[blue!70!black] (2.8,3.1) circle (2pt) node[right=3pt] {顺序 $B$};
  \fill[blue!70!black] (3.6,1.75) circle (2pt) node[right=3pt] {顺序 $C$（找到的最优）};
  \draw[-{Stealth}, gray] (2.15,3.8) to[bend left=20]
    node[midway, right=1pt] {\scriptsize 胜过} (2.75,3.25);
  \draw[-{Stealth}, gray] (2.95,3.0) to[bend left=20]
    node[midway, right=1pt] {\scriptsize 胜过} (3.55,1.9);
  \draw[decorate, decoration={brace, amplitude=4pt}] (5.9,1.75) -- (5.9,1.02)
    node[midway, right=6pt, align=left] {\scriptsize 认证\\[-2pt]\scriptsize 间隙};
\end{tikzpicture}
\caption{收缩顺序的相对评判 vs.\ 绝对评判。启发式可以把找到的顺序互相排名（上方），但只有
认证下界（下方）才能把找到的最优解变成带可证最优性差距的解。在我们的校准算例上，宽度下界与
显式 sweep \emph{相遇}：宽度间隙为零。}
\label{fig:yardstick}
\end{figure}

缺失的是一把\emph{绝对标尺}：一个可信的界 $L$，连同一个证明——\emph{任何}收缩顺序的代价都
不低于 $L$。有了这样的下界，任何找到的顺序 $S$ 都获得一张证书：$\mathrm{cost}(S)/L$ 界定了
它的最优性差距——正如 LP 对偶界为分支定界的当前解作证，或 UNSAT 证书为 SAT 搜索收尾。绝对
参照的碎片是存在的：穷举动态规划能为小网络找到保证最优的收缩序~\cite{pfeifer2014faster}，
精确 treewidth 求解器也会给其宽度界附上紧凑证书~\cite{tamaki2022heuristic}——但前者远在
要紧的实例尺寸之下就停了，且两者的保证都落在未经验证的搜索与检查代码上。TNCO 此前不存在的，
是一个其\emph{验证}本身就是机器核验定理的下界。

为什么没有？因为一个\emph{可信的}代价下界必须挡住三种相互独立的攻击，而每种攻击在历史上都
是致命的（图~\ref{fig:floors}）。第一，支撑下界的图论宽度界可能干脆是错的：宽度理论中的下
界证明微妙冗长，未经验证的界什么也认证不了。第二，宽度可能翻译不成代价：宽度数的是最大中间
张量的指标个数，而如果该张量具有低秩结构，更聪明的表示可以付得更少——decision diagram 模拟
器在结构化实例上确实击穿了宽度代理~\cite{dd2025treewidth}。第三，这个翻译必须对\emph{手头
的实例}成立，而不只在最坏情形下成立：浅层随机电路尽管宽度很大却是经典易模拟
的~\cite{napp2022efficient}。我们的回应是：让每一环都成为一条具名的机器核验定理，并且
明确分工——找一个好 bramble（下文的对偶证书）是启发式搜索问题，而\emph{验证}一个
bramble 是 verified 定理。前两种攻击对 stem 序被无条件封死；第三种由一张独立的框架层面证书
回应，它面向割面处的 \emph{generic} 参数族，并刻意不覆盖离散门集上的具体电路
（\S\ref{sec:endtoend}）。

\begin{figure}[t]
\centering
\resizebox{0.92\textwidth}{!}{%
\begin{tikzpicture}[font=\small,
  atk/.style={-{Stealth}, thick, red!60!black}]
  \node[draw, rounded corners, fill=blue!8, minimum width=6.4cm, align=center] (w) at (0,2.7)
    {\textbf{宽度下界} --- bramble 作为对偶证书\\[-1pt]
     {\scriptsize\ttfamily grid\_pathwidth\_exact}};
  \node[draw, rounded corners, fill=orange!12, minimum width=6.4cm, align=center] (c) at (0,1.35)
    {\textbf{代价下界} --- $\mathrm{cost}(o)=\sum_t 2^{|B_t|}$\\[-1pt]
     {\scriptsize\ttfamily gridModel\_cost\_floor}};
  \node[draw, rounded corners, fill=green!10, minimum width=6.4cm, align=center] (r) at (0,0)
    {\textbf{压缩下界} --- Contraction Rigidity\\[-1pt]
     {\scriptsize 仅覆盖 generic 参数族}\\[-1pt]
     {\scriptsize\ttfamily ContractionRigid.no\_compression}};
  \draw[atk] (-5.6,2.7) node[left, align=right, font=\scriptsize]
    {未经验证的\\图论界} -- (w.west);
  \draw[atk] (-5.6,1.35) node[left, align=right, font=\scriptsize]
    {宽度只是\\代理量} -- (c.west);
  \draw[atk] (-5.6,0) node[left, align=right, font=\scriptsize]
    {低秩逃逸\\（DD/MPS）} -- (r.west);
  \foreach \y in {2.7,1.35,0} {\draw[red!60!black, thick]
    (-4.55,\y-0.14) -- (-4.31,\y+0.14) (-4.55,\y+0.14) -- (-4.31,\y-0.14);}
\end{tikzpicture}}
\caption{一个可信的代价下界用三层\emph{分别}认证的结构挡住三种攻击；每一层都是一条仅依赖标
准公理的具名 Lean 定理。三层由 prose 而非单一定理关联（压缩下界是位于主定理合取之外的伴随
定理）：压缩下界谈论的是割面处的参数化张量
族，不是代价模型的边界张量；它在基准上的实例化是与实例无关的能力见证；且它只对
\emph{generic} 参数族挡住低秩攻击——离散门集上的具体电路（包括物理基准电路）在其适用范围
之外（\S\ref{sec:endtoend}、\S\ref{sec:limits}）。}
\label{fig:floors}
\end{figure}

本文为晶格校准族交付这把标尺，并附带一个 verified 证书验证接口——只要为任意网络供给一个
bramble，接口即适用（\S\ref{sec:limits}）——在 Lean~4 +
Mathlib~\cite{moura2021lean,mathlib2020} 中形式化：约
2{,}800 行、18 个模块（17 个证明模块加一个构建期公理审计模块）、无 \texttt{sorry}、零自定
义公理——下文每条定理的
\texttt{\#print axioms} 都只报告 \texttt{propext}、\texttt{Classical.choice}、
\texttt{Quot.sound}。无条件的下界量化在所有 \emph{sweep 型（stem）}顺序上——这正是
state-vector 类模拟器与 stem 型 supremacy 策略实际占据的类~\cite{huang2020classical}。推广
到任意收缩树依赖唯一一条显式陈述的已发表假设（晶格上 $\pw=\Theta(\tw)$）；开发中不存在树收
缩序的类型，因此该条件性陈述是 prose 中援引的文献推论，不是一条 Lean 定理
（\S\ref{sec:limits}）。

\paragraph{贡献。}
\begin{itemize}
  \item \textbf{TNCO 的绝对标尺（\S\ref{sec:framework}、\S\ref{sec:model}）。}一个形式化
    的代价/宽度模型，其中"没有任何 stem（sweep）收缩顺序低于界 $L$"是一条量化在\emph{真
    实}序空间上的定理
    ——序空间即网络图的真实路径分解，宽度全部派生、绝无设定——并配有显式代价对象：
    \texttt{orderCost} 对物化的边界张量求和，下界以\emph{代价}陈述读出
    （\texttt{gridModel\_cost\_floor}：$N\times N$ 晶格的每个序 $o$ 满足
    $2^{N+1}\le\mathrm{cost}(o)$——宽度下界的直接推论，显式化以使标尺为代价定价）。
  \item \textbf{Lean 中的 bramble 证书（\S\ref{sec:bramble}）。}bramble 是宽度的组合
    \emph{对偶证书}，扮演 LP 对偶之于分支定界下界的角色：无论 bramble 是如何找到的——手
    工、启发式搜索——被核验的定理都把它变成绝对下界。stem（sweep）序恰好是路径分解，其对偶
    论证只需\emph{区间} Helly——初等，不同于 treewidth 所需而 Mathlib 尚缺的子树 Helly。我
    们对真正的\emph{相触（touching）}bramble（块相交\emph{或}相邻）证明了该界，并实例化两
    个：初等十字（阶 $\Theta(k)$）与增广的 Seymour--Thomas 式 bramble
    （\S\ref{sec:exact}），后者的阶 $k+2$ 把 grid 下界钉到\emph{精确}。据我们所知，这是
    Lean/Mathlib 中首个 pathwidth 下界理论。
  \item \textbf{Contraction Rigidity：一个刻意初等的框架层面无压缩证书
    （\S\ref{sec:rigidity}）。}只有中间张量不发生秩坍缩，宽度下界才为代价定价。我们用一个
    独立、自立的证书回应这个漏洞：\texttt{BulkEntangling} $\Rightarrow$
    \texttt{ContractionRigid}——一
    个满秩配置使割面行列式成为非零多项式，秩坍缩被限制在该多项式的零集内，全程无概率。引擎
    刻意保持初等，证书的逻辑内容同样初等：求值与行列式交换，故 \texttt{ContractionRigid} 等
    价于行列式多项式非零这一条（\S\ref{sec:rigidity}）；价值在于封装：前件经 Kronecker 路由
    约化为组合的跨割面匹配，落点是一条形式化的\emph{无压缩定理}
    （\texttt{ContractionRigid.no\_compression}）：对刚性参数族，在一个非零多项式的零集之
    外，割面 flattening 不存在
    内维低于 $2^k$ 的因子分解——没有 MPS 式的 bond、没有任何秩揭示压缩
    （decision diagram 漏洞）。我们把这一层排在宽度下界之下——其数学内容是一个单点求值
    论证——其适用范围有三重边界。该零集的零测度解读只在 prose 中陈述
    （\S\ref{sec:limits}）；该定理\emph{没有}与 \S\ref{sec:model} 的代价下界复合（两张证
    书作为两条独立定理分别证明、互不复合，\S\ref{sec:endtoend}）；且证书约束的是\emph{给定的}族 $T$——
    它在基准割面上的存在式实例化由一个在任何平衡割面上都可得的 Bell 型见证兑现，不携带任
    何架构信息（\S\ref{sec:endtoend}）。
  \item \textbf{宽度上零间隙的认证校准算例（\S\ref{sec:endtoend}）。}在 $53\times53$ 深窗口时空
    grid——Sycamore-53 \emph{割面尺度}上的校准算例；\S\ref{sec:endtoend} 精确说明它与
    $(53,20)$ 基准参数的关系——上，一条机器核验定理把最优值精确钉死：$\cc=54$
    （\texttt{sycamore53\_cc\_exact}），每个 stem 序总代价 $\ge 2^{54}$
    （\texttt{sycamore53\_cost\_floor}），且显式 sweep 达到宽度下界：一个携带机器核验
    \emph{最优性证明}的展示解，其 \texttt{\#print axioms} 只报告标准公理。零间隙是
    \emph{宽度}陈述：代价下界继承的是宽度下界的指数、而非 sweep 的步数，sweep 自身的总代
    价在 $53^2\cdot 2^{54}$ 量级（\S\ref{sec:endtoend}）。一条置于该合取之外的伴随定理记
    录秩侧：一个由跨割面匹配构造的显式见证族——门参数为常量、沿块对齐匹配取恒等
    bond——在 $53$ 比特割面上达到满秩，仅认证该割面的纠缠\emph{能力}
    （\texttt{sycamore53\_bond\_floor} 进而界住对该见证的任何内维 $<2^{53}$ 的穿割分解）。
  \item \textbf{几何路由的脚手架（\S\ref{sec:routing}）。}为匹配前件提供组合支撑：
    $b\cdot d$ 深度封顶是一条关于抽象门调度的计数定理，bond 可达性是链与芯片\emph{耦合图}
    中的显式有界长度 walk。电路动力学模型未被形式化，且 \S\ref{sec:endtoend} 的刚性见证并
    不消费这些定理——它们为物理路由一步划定范围，而非将其兑现。
\end{itemize}

%% ===================================================================
\section{背景}\label{sec:background}

\subsection{张量网络、顺序与宽度}
量子电路编译为张量网络：每个门一个张量，每条线一个共享指标。模拟电路即把网络收缩为标量，
\emph{收缩顺序}调度两两合并。顺序的代价由其最宽的中间张量主导：一个有 $w$ 个开放指标的中间
张量持有 $2^w$ 个分量。Markov 与 Shi~\cite{markov2008simulating} 把所有顺序上的最优宽度与
网络线图的 treewidth 等同起来，这钉住了\emph{收缩范式}的复杂度（含指标切片——它以内存换时
间，但不减少总功）。收缩范式的界带有两条界定全文适用范围的限定。其一，它不是无条件硬度论断：利
用\emph{取值}结构的表示——stabilizer 态、decision
diagram~\cite{dd2025treewidth}——在结构化实例上能胜过宽度，且浅层随机 2D 电路本身就是经典
易模拟的~\cite{napp2022efficient}。因此我们的下界生活在收缩范式之内，且
\S\ref{sec:rigidity} 证明对 \emph{generic} 张量取值，"取值结构"这条逃逸通道对割面处的刚性
参数族也被封死（一张独立的伴随证书，从不与代价下界复合）。其
二，TNCO 的 NP-hardness 是关于一般网络的陈述；在平面 grid 上宽度是教科书已知的——而这恰恰
是晶格族在 \S\ref{sec:endtoend} 中够格充当\emph{校准}类的原因。

\subsection{stem 序即路径分解}
在 supremacy 型电路的实践中占主导的顺序是 \emph{stem}~\cite{huang2020classical}：一个不断
长大的张量把网络一块块吸收进来，恰如 state vector 扫过电路。用树的语言说，收缩树是一条
caterpillar——带叶子的脊柱。本文核心的组合恒等式（图~\ref{fig:stemsweep}）是：
\textbf{一个 stem 序与网络图的一个路径分解是同一份数据}：第 $t$ 个 bag $B_t$ 是吸收 $t$ 步
后已吸收区域的边界，序的宽度是最大 bag，最优 stem 宽度即 pathwidth。正是这个恒等式使 stem
序的认证下界在今天即可形式化：pathwidth 的下界依赖\emph{区间} Helly（\S\ref{sec:bramble}），
一个一维论证；而 treewidth 需要 Mathlib 尚不具备的子树 Helly 性质。

\begin{figure}[t]
\centering
\begin{tikzpicture}[scale=0.52, font=\scriptsize]
  \begin{scope}
    \fill[gray!15] (-0.4,-0.4) rectangle (1.45,4.4);
    \foreach \x in {0,...,4} \foreach \y in {0,...,4} {
      \draw[gray] (\x,\y) circle (1.6pt);
      \ifnum\x<4 \draw[gray!70] (\x+0.08,\y) -- (\x+0.92,\y); \fi
      \ifnum\y<4 \draw[gray!70] (\x,\y+0.08) -- (\x,\y+0.92); \fi
    }
    \foreach \y in {0,...,4} \fill[orange] (2,\y) circle (2.6pt);
    \draw[very thick, dashed, orange!80!black] (1.55,-0.5) -- (1.55,4.5);
    \node[orange!80!black] at (2.0,5.05) {bag $B_t$ = 前沿};
    \node[gray!60!black] at (0.5,-1.0) {已吸收};
    \draw[-{Stealth}, thick] (0.3,4.9) -- (1.3,4.9) node[midway, above] {扫描};
  \end{scope}
  \node at (6.3,2) {\Large$\;=\;$};
  \node[align=center] at (6.3,2.8) {同一份\\数据};
  \begin{scope}[shift={(8.2,0)}]
    \foreach \i in {0,...,5} \fill[blue!70!black] (\i*1.1,2.6) circle (2.6pt);
    \foreach \i in {0,...,4} \draw[blue!70!black, very thick]
      (\i*1.1+0.07,2.6) -- (\i*1.1+1.03,2.6);
    \foreach \i/\n in {0/1,1/2,2/1,3/2,4/1,5/1} {
      \ifnum\n=1
        \draw[gray] (\i*1.1,2.52) -- (\i*1.1,1.7); \draw[gray] (\i*1.1,1.7) circle (1.6pt);
      \else
        \draw[gray] (\i*1.1,2.52) -- (\i*1.1-0.3,1.7); \draw[gray] (\i*1.1-0.3,1.7) circle (1.6pt);
        \draw[gray] (\i*1.1,2.52) -- (\i*1.1+0.3,1.7); \draw[gray] (\i*1.1+0.3,1.7) circle (1.6pt);
      \fi
    }
    \node[blue!70!black] at (2.75,3.25) {脊柱（不断长大的张量）};
    \node[gray!60!black] at (2.75,1.1) {已吸收的叶子};
  \end{scope}
\end{tikzpicture}
\caption{stem（caterpillar）收缩序\emph{就是}路径分解：脊柱上第 $t$ 个中间张量的开放指标恰
为前沿 bag $B_t$。序的宽度 $=$ 最大 bag；最优 stem 宽度 $=$ pathwidth。}
\label{fig:stemsweep}
\end{figure}

\subsection{扫描割面与基准尺度}
对 $n$ 比特、深度 $d$ 的晶格电路，两个典范扫描给出割面尺寸 $n$（时间向）与
$\lfloor\sqrt{n}\rfloor\cdot d$（空间向），故可达的宽度尺度为
$\mathrm{sweepCut}(n,d)=\min(n,\lfloor\sqrt{n}\rfloor d)$；在 Sycamore-53 参数
$(n,d)=(53,20)$~\cite{arute2019quantum} 下这个值是 $53$，而民间流传的代价估计
$4\cdot 2^{53}\cdot 150\approx 10^{18}$——其中步数 $150$ 取自该民间估计的输入，并非推导
所得——是纯算术，我们的开发同样将其认证
（\texttt{CorollaryD}）。全文中"Sycamore-53 尺度"指这组参数固定的割面值；我们实例化的晶
格是线链的 $53\times53$ \emph{深窗口}时空 grid，其深度窗口（$53$ 层）超过基准深度
（$20$）——\S\ref{sec:endtoend} 说明每个认证合取项谈论的是哪个对象。

\subsection{Lean 4 与 Mathlib}
开发构建于 Lean~4 与 Mathlib（均钉在 v4.30.0），使用 \texttt{SimpleGraph} 盒积表示晶格、
\texttt{Finset} 组合学表示 bramble 与 transversal、\texttt{MvPolynomial} 加矩阵行列式驱动
genericity 引擎。审计接口是 \texttt{\#print axioms}：本文引用的每条定理都报告
\texttt{[propext, Classical.choice, Quot.sound]}——库中不存在任何自定义公理，我们依赖的唯
一文献输入（晶格上 $\pw=\Theta(\tw)$）以显式假设出现在定理陈述\emph{内部}，从不作为环境假定。

%% ===================================================================
\section{形式化框架}\label{sec:framework}

整个开发干净地分成线性代数与组合两半，而让两半都保持简单的设计决策是：互不引入对方的机制。
线性代数一半围绕一个接口组织——\emph{张量跨割面的 flattening}——组合一半从不需要它；两半
只在 \S\ref{sec:endtoend} 的合取中相遇。

线性代数一侧，腿集为 $I$ 的张量就是每条腿一个比特所索引的数组，割面是腿上的谓词 $p$；
flattening 把张量重排成矩阵，行是 $p$ 侧的赋值，列是其余腿的赋值：

\begin{lstlisting}
abbrev QTensor (I R : Type*) := (I → Fin 2) → R

def flatten (T : QTensor I R) :
    Matrix ({i // p i} → Fin 2) ({i // ¬ p i} → Fin 2) R :=
  Matrix.of fun a b => T ((cutEquiv p).symm (a, b))
\end{lstlisting}

该矩阵的秩即跨割面的 Schmidt 秩，而"无秩坍缩"——\S\ref{sec:rigidity} 所认证的性质——恰是
平衡割面处其方形形式的非奇异性。一条引理承担全部输运工作：逐元求值与 flattening 交换
（\texttt{flatten\_map}），因此对多项式元证得的任何 genericity 陈述都下降到每个具体取值。

组合一侧，电路图由其收缩序空间建模：类型 \texttt{Order}、每个序的派生量 \texttt{width}、谓
词 \texttt{IsStem}，以及两个最优值 $\cc$（宽度在全体序上的下确界）与 $\pw$（在 stem 序
上）。已证的骨架小而无条件：$\cc\le\pw$、stem 最优被达到、以及\emph{给定}界
$\pw\le C\cdot\cc$ 时存在与最优相差因子 $C$ 之内的 stem 序。文献输入被刻意做成\emph{非}公
理：\texttt{MarkovShi G} 与 \texttt{LatticePathwidthBound G C} 是具名 \texttt{Prop}，调用
者必须为手头的具体模型供给。在 \S\ref{sec:model} 的模型上后者平凡成立——该模型的序空间只
含 stem 序，此界是同一个下确界与自身相比——因此本文没有任何定理携带未证的环境假设；真正的
树序内容恰是该模型无法表达的部分（\S\ref{sec:limits}）。

%% ===================================================================
\section{Contraction Rigidity}\label{sec:rigidity}

只有最宽的中间张量真正携带 $2^w$ 个自由度，宽度下界才认证代价。文献把这当作 average-case
问题——"Haar 随机电路的割面以高概率满秩"——并用随机矩阵分析作答，而后者的保证对 regime
敏感：浅层随机 2D 电路本身就是经典易模拟的~\cite{napp2022efficient}。我们把概率问题换成
结构问题，答案变成单点多项式非零性的两行推论。

称参数化张量族 $T$（元是门参数的多项式）在平衡割面处 \texttt{BulkEntangling}，若\emph{某
个}参数点使割面 flattening 非奇异——即架构\emph{有能力}达到满 Schmidt 秩。称它
\texttt{ContractionRigid}，若 flattening 的行列式是非零多项式，且在该多项式零集之外的每个
参数处 flattening 非奇异。结构定理是：能力升级为 genericity：

\begin{lstlisting}
theorem contraction_rigidity
    (e : ({i // p i} → Fin 2) ≃ ({i // ¬ p i} → Fin 2))
    {T : QTensor I (MvPolynomial ι K)}
    (h : BulkEntangling p e T) : ContractionRigid p e T
\end{lstlisting}

引擎是初等的——在一个取值处非零的行列式是非零多项式——且有一件簿记事实应当明说。由于求值
与行列式交换（\texttt{det\_squareFlatten\_map}，一条无条件恒等式），\texttt{ContractionRigid}
的第二子句（零集之外逐点非奇异）在第一子句成立时对\emph{一切}族自动成立，因此
\texttt{ContractionRigid} 逻辑等价于其第一子句本身：flattening 的行列式是非零多项式。定义保
留逐点子句，只因下游消费者（下文的 \texttt{no\_compression}）恰以该形态使用它。零集的
properness（补集非空，需无穷域）另行证明
（\texttt{exists\_nonsingular\_of\_det\_ne\_zero}），不属于该证书。不含概率；零测度推论（对
任意连续门测度 $\Pr[\text{秩坍缩}]=0$）是一步之遥的推论，留作后续工作。

让该定理可用的关键在于其前件是\emph{组合的}（图~\ref{fig:matching}）。若割面两侧各 $k$ 个
比特能被一个完美匹配配对、每对携带一个非奇异的两比特 bond，则见证配置下的 flattening 是各
bond 的 Kronecker 积，而非奇异块的积非奇异（\texttt{det\_bondProd\_ne\_zero}，对 $k$ 归纳
加 \texttt{det\_kronecker}）。于是为一个架构检验 \texttt{BulkEntangling} 约化为展示一个跨
割面匹配——一个图条件，其硬件图一侧由 \S\ref{sec:routing} 划定范围（并不兑现物理路由这一
步）。该前件对真实门集非空：开发
证明了标准门库（$X, Y, Z, H, S, T, \sqrt{X}, \mathrm{CZ}, \mathrm{SWAP}, \mathrm{CNOT},
\mathrm{iSWAP}, \mathrm{fSim}(\theta,\phi)$ 等）全部 $\det\neq 0$，因此恒等 bond、Bell 对
和 Sycamore 原生的 $\mathrm{fSim}$ 都合格。必须划清一条边界：定理约束的是\emph{给定的}族
$T$（某个固定电路的张量）；相比之下，\S\ref{sec:endtoend} 伴随能力定理所用的裸存在式
$\exists T$ 在\emph{任何}平衡割面上都由 Bell 见证兑现，因此它只认证割面尺寸层面的能力，从
不构成对特定架构的性质断言。

\begin{figure}[t]
\centering
\begin{tikzpicture}[scale=0.75, font=\small]
  \draw[dashed, thick] (2.1,-0.6) -- (2.1,3.6) node[above] {割面};
  \foreach \y/\j in {0/4,1/3,2/2,3/1} {
    \draw[fill=gray!15] (0.9,\y) circle (7pt) node {$L_{\j}$};
    \draw[fill=gray!15] (3.3,\y) circle (7pt) node {$R_{\j}$};
    \draw[blue!70!black, very thick] (1.16,\y) -- (3.04,\y)
      node[midway, above] {\scriptsize $B_{\j}$};
  }
  \node[align=left, anchor=west] at (4.6,2.4)
    {每个 bond：$\det B_j \neq 0$};
  \node[align=left, anchor=west] at (4.6,1.6)
    {$\Rightarrow$ flattening $=\bigotimes_j B_j$};
  \node[align=left, anchor=west] at (4.6,0.8)
    {$\Rightarrow$ 在\emph{一个}点达到 Schmidt 秩 $2^{k}$};
  \node[align=left, anchor=west] at (4.6,0.0)
    {$\Rightarrow$ \emph{generic} 满秩};
\end{tikzpicture}
\caption{Kronecker 路由把刚性前件约化为组合：带非奇异 bond 的跨割面完美匹配使割面
flattening 成为 Kronecker 积，于是在一个门配置处满秩——而单点求值论证把这一个点升级为
全体 generic 参数。全程无概率。}
\label{fig:matching}
\end{figure}

\paragraph{从秩到压缩：无压缩定理。}
最后一步是刚性对代价叙事的一个自立补充。若割面 flattening $M$ 经内维为 $r$ 的指标集分解为
$M=A\cdot B$——一个 MPS 式的 bond、一个 decision diagram 的节点割、任何秩揭示压缩——则
$r$ 从上方界住 $M$ 的秩。形式化地，对刚性族及零集之外的任何参数：

\begin{lstlisting}
theorem ContractionRigid.no_compression ... :
    squareFlatten p e (fun x => eval v (T x)) = A * B →
    2 ^ Fintype.card {i // p i} ≤ r
\end{lstlisting}

\noindent （经 \texttt{M.mulVec} 的单射性与 \texttt{finrank} 在单射下的单调性证得）。因此
对刚性参数族，在一个非零多项式的零集之外，\emph{没有任何}表示能以低于 $2^k$ 的内维穿过割面
——该零集对任意连续门测度是零测的，这个解读我们陈述但未形式化（\S\ref{sec:limits}）。结构
化实例——stabilizer 电路、真正低秩的张量——恰好
生活在那个零集内并且可以击穿下界；这正是 generic 实例证书的既定语义。范围说明一
条：\texttt{no\_compression} 谈论的是割面处的参数化张量族；没有任何定理把它与
\S\ref{sec:model} 的 \texttt{orderCost} 联结——压缩下界与代价下界在
\S\ref{sec:endtoend} 中作为两条\emph{独立}定理分别认证，互不复合。

%% ===================================================================
\section{经由区间 Helly 的 stem 宽度}\label{sec:bramble}

宽度下界是 pathwidth 下界，其形式化接口刻意保持极小。一个候选 stem 序是一个
\texttt{LinearBags}——有限的 bag 序列——连同路径分解的三条定义公理：每个顶点所在的 bag 构
成连续区间（\texttt{vint}）、每条边落在某个 bag 内（\texttt{edge}）、每个顶点被覆盖
（\texttt{cov}）。下界所需的一切都从这三条导出。

对偶证书是 \emph{bramble}：一族连通、两两相触的顶点块。其\emph{阶}是 transversal（与每个块
相交的集合）的最小尺寸，而要证的界是：任何路径分解的最大 bag 至少是该阶。证明
（图~\ref{fig:helly}）正是 pathwidth 限制开始回报的地方。连通性加上边覆盖公理迫使每个块沿
一段连续的指标\emph{区间}与 bag 相遇（\texttt{bag\_meets\_betweenness}，对 walk 归纳）；两
两相触迫使这些区间两两重叠；而两两重叠的区间有公共点——区间 Helly 性质，其全部证明就是"取
最大的左端点"（\texttt{interval\_helly}，六行）。公共指标处的 bag 与每个块相遇，因此是一个
transversal，其尺寸界住阶。对 treewidth 同样的论证需要子树的 Helly 性质，Mathlib 没有；对
pathwidth 一维情形即够，这使认证 stem 下界落在当前 Mathlib 的能力范围之内。

\begin{figure}[t]
\centering
\begin{tikzpicture}[scale=0.72, font=\scriptsize]
  \foreach \i in {1,...,8} {
    \draw[gray!60, fill=gray!8] (\i-0.42,0) rectangle (\i+0.42,0.7);
    \node at (\i,0.35) {$B_{\i}$};
  }
  \draw[blue!70!black,   line width=3pt] (1.6,1.15) -- (6.4,1.15)
    node[pos=0.1, above] {块 $\beta_1$};
  \draw[orange!85!black, line width=3pt] (3.6,1.6)  -- (8.4,1.6)
    node[pos=0.9, above] {块 $\beta_2$};
  \draw[green!55!black,  line width=3pt] (0.6,2.05) -- (5.4,2.05)
    node[pos=0.1, above] {块 $\beta_3$};
  \draw[dashed, thick, red!70!black] (5,-0.35) -- (5,2.5);
  \node[red!70!black, align=center] at (5,2.9)
    {公共 bag（区间 Helly）：\\与\emph{每个}块相遇 $\Rightarrow$ $|B|\ge$ 阶};
\end{tikzpicture}
\caption{pathwidth 的对偶论证。连通块沿路径分解的连续指标\emph{区间}与 bag 相遇；两两相触
的块给出两两重叠的区间；区间 Helly（最大左端点）产出一个与所有块相遇的 bag，因此最大 bag
至少是 bramble 的 transversal 阶。只需要一维 Helly，这使 stem 序下界在当前 Mathlib 中即可形式化。}
\label{fig:helly}
\end{figure}

两个改进使该界超越玩具 bramble 可用。第一，块可以\emph{相触而不相交}——真正的 bramble 条件
允许相邻——此时公共指标来自边覆盖公理而非公共顶点
（\texttt{pathwidth\_ge\_order\_of\_touching}）。第二，整个开发以任意有限类型为索引，因此
结构化的 bramble（\S\ref{sec:exact}）可以自然编码。用初等的\emph{十字} bramble——第 $i$ 块
是第 $i$ 行并第 $i$ 列，经盒积 walk 嵌入证明其在真实 grid 图中连通——实例化，即得到完全无
条件的 $\Theta(k)$ 下界：

\begin{lstlisting}
theorem grid_pathwidth_lower_unconditional (k : ℕ) (hk : 0 < k)
    (L : LinearBags (Fin k × Fin k)) (hlen : 0 < L.len)
    (hvint : ...) (hedge : ...) (hcov : ...) :
    k ≤ 2 * L.maxBag hlen
\end{lstlisting}

%% ===================================================================
\section{精确下界：增广 bramble}\label{sec:exact}

对渐近学而言，$\Theta(k)$ 就是故事的终点；对一把\emph{标尺}而言不是。用 $k/2$ 的下界对上
$k+1$ 的 sweep，只能认证"指数意义下相差不超过两倍"——一张弱证书，而且是不必要的弱，因为松
掉的部分全在 bramble 里。十字 bramble 的阶真的只有 $\lceil k/2\rceil$：一个顶点 $(i,j)$ 同
时命中两个十字。要关掉这个间隙需要更好的对偶见证——Seymour--Thomas 风
格~\cite{seymour1993graph}的增广 bramble（图~\ref{fig:augmented}）：在
$(k{+}1)\times(k{+}1)$ 的 grid 上，取 $k\times k$ 子网格的 $k^2$ 个十字，加上末行，加上去
掉角点的末列。

\begin{figure}[t]
\centering
\begin{tikzpicture}[scale=0.62, font=\scriptsize]
  \fill[blue!6] (-0.45,0.55) rectangle (4.45,5.45);
  \foreach \x in {0,...,5} \foreach \y in {0,...,5} {
    \ifnum\x<5 \draw[gray!50] (\x+0.1,\y) -- (\x+0.9,\y); \fi
    \ifnum\y<5 \draw[gray!50] (\x,\y+0.1) -- (\x,\y+0.9); \fi
  }
  \foreach \x in {0,...,5} \foreach \y in {0,...,5} \draw[gray] (\x,\y) circle (1.7pt);
  \foreach \x in {0,...,4} \fill[blue!70!black] (\x,3) circle (2.6pt);
  \foreach \y in {1,...,5} \fill[blue!70!black] (2,\y) circle (2.6pt);
  \node[blue!70!black] at (0.85,4.6) {十字 $(i,j)$};
  \foreach \x in {0,...,5} \fill[red!75!black] (\x,0) circle (2.6pt);
  \node[red!75!black] at (6.35,0) {末行};
  \foreach \y in {1,...,5} \fill[green!55!black] (5,\y) circle (2.6pt);
  \node[green!55!black] at (6.55,3.4) {末列};
  \draw[very thick, red!75!black]   (2,0.15) -- (2,0.85);
  \draw[very thick, green!55!black] (4.15,3) -- (4.85,3);
  \node[align=left, anchor=west] at (-0.2,-1.2)
    {仅相触的边：十字与边界块\emph{不相交}};
\end{tikzpicture}
\caption{$(k{+}1)\times(k{+}1)$ grid 上的增广 bramble：$k^2$ 个子网格十字（蓝）、末行
（红）、去角末列（绿）。边界块只通过边与十字相触，恰由相触 bramble 界
（\texttt{pathwidth\_ge\_order\_of\_touching}）经边覆盖公理处理。计数：transversal 需要一
个红顶点、一个绿顶点、以及 $k$ 个子网格顶点——否则某行 $i_0$ \emph{且}某列 $j_0$ 同时被漏
掉，十字 $(i_0,j_0)$ 无人命中——故阶 $\ge k+2$，与 sweep 精确相遇。}
\label{fig:augmented}
\end{figure}

边界块与每个十字\emph{不相交}——它们只通过边相触——这正是
\S\ref{sec:bramble} 的相触 bramble 界为之而生的场景。子网格十字的连通性需要不越出子
网格的 walk；我们构造带显式支撑界的递增链 walk，并把两条臂都路由到十字中心。而
transversal 计数是经典的"漏一行且漏一列"论证，用投影形式化：transversal 必须包含一个末行顶
点与一个末列顶点（两个不属于任何十字的不同顶点），且其子网格部分必须覆盖全部 $k$ 个子网格
行或全部 $k$ 个子网格列——否则某个十字 $(i_0,j_0)$ 无人命中——故阶至少 $k+2$：

\begin{lstlisting}
theorem grid_pathwidth_exact (k : ℕ) (hk : 0 < k)
    (L : LinearBags (Fin (k+1) × Fin (k+1))) (hlen : 0 < L.len)
    (hvint : ...) (hedge : ...) (hcov : ...) :
    k + 2 ≤ L.maxBag hlen
\end{lstlisting}

\noindent 这与 \S\ref{sec:model} 的滑动窗口 sweep 精确相遇：grid 的 stem 最优被精确钉死，
而不只是渐近地夹在区间里。

%% ===================================================================
\section{忠实的电路模型}\label{sec:model}

一个下界的意义取决于它量化在哪个序空间上，而形式化最脆弱的地方恰在于此：对一个自己设定数值
的模型证明强命题是容易的。我们的第一版模型正是如此——单个抽象序、恒定宽度——我们汇报这条
死路，因为对它的修复正是方法论要点：在最终模型中\textbf{没有任何宽度或代价数值是设定的}
（记录类型中保留两个惰性字段 \texttt{tw} 与 \texttt{IsLattice} 以兼容接口，本文引用的定理
均不消费它们）。一个序\emph{就
是}真实 grid 图的路径分解（\S\ref{sec:bramble} 的三条公理，打包为 \texttt{IsPathDecomp}），
其宽度\emph{就是}它的最大 bag，其代价\emph{就是}它物化的边界张量总尺寸，
$\mathrm{cost}(o)=\sum_t 2^{|B_t|}$（\texttt{orderCost}）。

\begin{figure}[t]
\centering
\begin{tikzpicture}[scale=0.62, font=\tiny]
  \foreach \x in {0,...,4} \foreach \y in {0,...,4} {
    \ifnum\x<4 \draw[gray!40] (\x+0.12,4-\y) -- (\x+0.88,4-\y); \fi
    \ifnum\y<4 \draw[gray!40] (\x,4-\y-0.88) -- (\x,4-\y-0.12); \fi
  }
  \foreach \x in {0,...,4} \foreach \y in {0,...,4} {
    \pgfmathtruncatemacro{\p}{\x*5+\y}
    \ifnum\p<8
      \draw[fill=gray!45] (\x,4-\y) circle (5pt); \node at (\x,4-\y) {\p};
    \else\ifnum\p<14
      \draw[line width=1.4pt, orange!90!black, fill=orange!12] (\x,4-\y) circle (5pt);
      \node at (\x,4-\y) {\p};
    \else
      \draw (\x,4-\y) circle (5pt); \node at (\x,4-\y) {\p};
    \fi\fi
  }
  \node[align=left, anchor=west, font=\scriptsize] at (5.0,3.4)
    {\textcolor{gray}{$\bullet$} 已吸收（$\mathrm{pos}<t$）};
  \node[align=left, anchor=west, font=\scriptsize] at (5.0,2.6)
    {\textcolor{orange!90!black}{$\circ$} bag $=$ 窗口 $[t,\,t{+}k]$};
  \node[align=left, anchor=west, font=\scriptsize] at (5.0,1.8)
    {$\circ$ 未来};
  \node[align=left, anchor=west, font=\scriptsize] at (5.0,1.0)
    {$|B_t| \le k{+}1$（位置单射）};
\end{tikzpicture}
\caption{滑动窗口 sweep（\texttt{sweepBags}，图示 $5\times5$ grid、$t=8$）：顶点按列主序位
置吸收；第 $t$ 步的 bag 是位置窗口 $[t, t+k]$——活跃前沿加上正被吸收的顶点。每个窗口至多
$k+1$ 个顶点，实现教科书 grid pathwidth，并与图~\ref{fig:augmented} 的下界相遇。}
\label{fig:window}
\end{figure}

\begin{lstlisting}
def gridModel (k : ℕ) : CircuitGraph where
  Order  := {L : LinearBags (Fin k × Fin k) // IsPathDecomp k L}
  width  := fun o => bagsWidth o.1
  IsStem := fun _ => True
  ...
\end{lstlisting}

\noindent 这里 \texttt{IsStem} $\equiv$ \texttt{True} 是定义性的而非退化的：该序空间恰由线
性（stem）序构成，而把量化扩展到树序正是 $\pw=\Theta(\tw)$ 文献输入，在用到之处以假设携带。

下界是 \S\ref{sec:exact} 的定理，如今直接谈论模型的每一个序。相配的上界是显式滑动窗口
sweep（图~\ref{fig:window}）：按列主序位置吸收顶点，令第 $t$ 步的 bag 为位置窗口
$[t, t+k]$——活跃前沿加正被吸收的顶点。窗口满足全部三条分解公理（区间性质是位置上的算术；
每条 grid 边的位置差为 $1$ 或 $k$，故较小端点处的窗口覆盖它），且由位置单射至多 $k+1$ 个顶
点。下界与 sweep 相遇，模型的最优被精确钉死，代价下界由
$2^{\mathrm{width}}\le\mathrm{cost}$ 推出：

\begin{lstlisting}
theorem gridModel_cc_eq (N : ℕ) (hN : 2 ≤ N) :
    (gridModel N).cc = N + 1

theorem gridModel_cost_floor (N : ℕ) (hN : 2 ≤ N)
    (o : (gridModel N).Order) : 2 ^ (N + 1) ≤ orderCost o
\end{lstlisting}

\noindent 最优更进一步是\emph{被达到的}（经 \texttt{Nat.sInf\_mem}），且
\texttt{LatticePathwidthBound (gridModel N) 1} 由构造成立——且是平凡地成立，因为该序空间
只含 stem（\S\ref{sec:framework}）——于是 \S\ref{sec:framework} 的每条假设式定理都无剩余
假设地适用于该模型，公理足迹恰好停在三条标准公理。

%% ===================================================================
\section{案例研究：宽度上零间隙的认证基准算例}\label{sec:endtoend}

为什么在晶格上校准？恰恰因为晶格\emph{不是}困难情形：其宽度是教科书已知的，所以它是唯一能
拿标尺读数与独立真值对表的实例族。把模型实例化在 $N=53$（53 线链的深窗口时空 grid），一条
定理合取整个证书。该算例与 Sycamore-53 基准的唯一联系，是下文合取项 (i) 固定的共享割面值
$53$；Lean 前缀 \texttt{sycamore53} 命名的是这个校准算例，不是物理器件：

\begin{lstlisting}
theorem sycamore53_lower_bound :
    CorollaryD.sweepCut 53 20 = 53 ∧
    ((10:ℕ)^18 ≤ CorollaryD.stemCost 53 20 150 ∧
      CorollaryD.stemCost 53 20 150 < (10:ℕ)^19) ∧
    (sycamore53.cc = 54 ∧
      (∃ o, sycamore53.IsStem o ∧
        sycamore53.width o = sycamore53.cc) ∧
      (∀ o : sycamore53.Order,
        2 ^ 54 ≤ GridModel.orderCost o))
\end{lstlisting}

作为证书来读，三个合取说的是：(i)~民间 $10^{18}$ 估计背后的扫描割面算术是精确的；(ii)~模型
的时空 grid 图的\emph{全体} stem 序上的最优值是 $54$——\S\ref{sec:exact} 的下界与
\S\ref{sec:model} 的 sweep 相遇，故该 sweep 是一个携带机器核验\emph{最优性证明}的展示解；
(iii)~每个序的总代价至少 $2^{54}$。对该定理执行
\texttt{\#print axioms} 报告 \texttt{[propext, Classical.choice, Quot.sound]}。

\paragraph{秩侧伴随定理。}秩侧被刻意排除在该定理的合取之外：它的逻辑内容比上述合取项单薄，
若并入主定理合取，便会借走后者的分量。一条伴随定理（\texttt{sycamore53\_bond\_floor}）展示
一个由跨割面匹配构造的见证张量族在 $53$ 比特割面上收缩刚性——对该族，
\texttt{no\_compression} 给出没有内维低于 $2^{53}$ 的因子分解穿过割面。这是关于割
面纠缠容量的存在性陈述，不是对任何具体基准电路的秩断言；该见证族在门参数上是
\emph{常量}（沿块对齐匹配取恒等 bond），其 \texttt{ContractionRigid} 的例外集为空。更强
地说：该存在式由一个在\emph{任何}指标类型的\emph{任何}平衡割面上都存在的 Bell 型见证兑
现，因此伴随定理除割面尺寸 $53$ 之外不携带该算例的任何特有信息——它认证的是该尺寸上的
\emph{能力}，generic 量词在这个见证上不做任何额外的功，实质性的刚性陈述仍是针对\emph{给
定}电路族的框架定理（\S\ref{sec:rigidity}）。

保真度范围。三个合取项与秩侧伴随定理谈论的并非同一个物理对象。合取项 (i)--(ii) 是基准参数
$(53,20)$ 处的算术，其割面值 $53$ 来自\emph{二维芯片}几何（每层
$\lfloor\sqrt{n}\rfloor$ 个跨割面门）。合取项 (iii) 认证的是\emph{线链}的 $53\times53$
时空 grid——一个 $53$ 层的深窗口，超过基准深度 $20$；在深度 $20$ 处线链的时空 grid 是
$53\times20$，其 stem 尺度由深度支配（$\approx 21$），且按 \S\ref{sec:routing} 的饱和判
据位于深 regime 阈值之下。(iii) 之下还有一步建模：被认证的图是完整的 $53\times53$ grid
图（两条路径的盒积），它是深度 $53$ 的 brickwork 链更稀疏的门级网络图的代理几何——两者在
$\Theta$ 尺度上一致、常数上可有出入，而定理量化的对象只是这个 grid 图。因此 (iii) 认证的
是 Sycamore-53 割面尺度上的深 regime
\emph{校准}算例——不是 $(53,20)$ 的线链，也不是物理芯片的 $(2{+}1)$ 维晶格。两个代价读数
同样是不同的度量、在数值上并不互洽：合取项 (ii) 以 bond 语义为 $150$ 个民间 stem 步定价
（$4\cdot 2^{53}\cdot 150\approx 5\times10^{18}$），而合取项 (iii) 的 \texttt{orderCost}
对模型的 $53^2=2809$ 个 bag 步各收 $2^{|B_t|}$ 个分量，显式 sweep 自身的
\texttt{orderCost} 约为 $53^2\cdot 2^{54}\approx 5\times10^{19}$——落在合取项 (ii) 的区
间之外。两个数字互不界定对方的对象；一种语义下的断言不得在另一种语义下解读。伴随能力定理
以 generic 形式陈述：具体的 Sycamore 电路从离散集合中取单比特门，真子簇陈述不覆盖它——
证书对那个具体实例不做秩断言。证书的匹配一侧已经覆盖二维芯片（\S\ref{sec:routing}）；芯
片晶格的宽度下界是后续工作（\S\ref{sec:limits}），且本节没有任何内容依赖它。

%% ===================================================================
\section{实证演示：拿标尺量一次真实搜索}\label{sec:empirical}

我们把标尺对着一次真实搜索做了检验。对一列直到基准尺度的 $N$，我们构
造封闭的 $N\times N$ grid 张量网络（所有 bond 维数为 $2$），运行 cotengra 的
hyper-optimizer~\cite{gray2021hyper}（v0.7.1，\texttt{minimize=size}，每实例 $128$ 次重
复，单 CPU 节点，总计约 22 分钟；对漂移实例另做 $512$ 与 $4096$ 次重复的预算敏感性重跑，见
下文），记录找到的最优收缩宽度——最大中间张量的 $\log_2$——并与
认证 stem 尺度 $N+1$ 对比（图~\ref{fig:gapplot}）。一条语义说明为对比划定范围：找到的顺序是
以 \emph{bond} 数计量的\emph{树}序，而认证下界是关于以\emph{顶点 bag} 计量的 \emph{stem}
序的陈述，两种计量只在相差一个加性单位的意义下一致——逐顶点推进的前沿多携带正被吸收顶点的
那一条 bond，而树式合并可以把它预先收缩掉。（具体地，在封闭 grid 上：列中位置的 sweep 中间
张量暴露 $N$ 条跨列 bond 加吸收点处的一条列内 bond，故其 bond 宽度为 $N{+}1$，与 bag 计数一
致；预收缩正被吸收的顶点恰好省下那一个单位。）因此跨两种语义做 $\pm 1$ 粒度的比较不在我们定理
的许可范围内；理论的预言是 $|\text{found}-(N{+}1)|\le 1+\text{搜索松弛}$。这个预言构成一次真正的检验，
因为在晶格上树序本就不该实质性胜过 stem：grid 上 $\pw=\tw$ 正是 stem 存在性论题的内容
（\S\ref{sec:limits} 显式携带的那条假设），所以一个可以自由返回\emph{任意}树的搜索，理应落
在 stem 线一个换算单位之内——而且绝不应实质性低于它。

\begin{figure}[t]
\centering
\begin{tikzpicture}[font=\scriptsize, xscale=0.145, yscale=0.62]
  \draw[-{Stealth}] (4,-1.9) -- (4,2.9) node[above left] {间隙};
  \draw (4,-1.9) -- (56,-1.9);
  \foreach \g/\l in {-1/-1, 0/0, 1/+1, 2/+2} {
    \draw (3.5,\g) -- (4,\g); \node[left] at (3.5,\g) {\l};
  }
  \foreach \n in {6,12,20,28,36,44,53} \node[below] at (\n,-1.9) {\n};
  \node[below=9pt] at (30,-1.9) {grid 尺寸 $N$};
  \draw[red!70!black, very thick] (4.5,0) -- (55.5,0);
  \node[red!70!black, above] at (16,0.10) {认证 stem 尺度 $N{+}1$};
  \foreach \n/\g in {6/-1, 8/0, 10/0, 12/0, 16/0, 20/0, 24/0, 28/0, 32/1,
                     36/2, 40/-1, 44/1, 48/0, 53/2}
    \fill[blue!70!black] (\n,\g) circle (2.4pt);
  \draw[gray, -{Stealth}] (44.5,-1.35) -- (40.6,-1.1);
  \node[gray, right, align=left] at (44.7,-1.35)
    {bag 对 bond 换算\\（$\pm 1$ 语义容差）};
  \draw[gray, -{Stealth}] (28.5,2.35) -- (35.3,2.08);
  \node[gray, left, align=right] at (28.3,2.35)
    {搜索方差——\\没有绝对参照\\就不可见};
\end{tikzpicture}
\caption{找到的收缩宽度减认证 stem 尺度 $N+1$，遍历 grid 族（cotengra hyper-optimizer，每
实例 $128$ 次重复，\emph{不受限的树}序）。所有点都落在 $[-1,+2]$ 内：十四个中八个与认证尺
度数值重合，两个位于 $-1$（在 bag 对 bond 的一个换算单位容差之内——并未低于 stem 论题），
四个漂到 $+1$/$+2$ 之上——一种相对基准测试无法察觉、而标尺一眼读出的质量缺口。}
\label{fig:gapplot}
\end{figure}

数据落在三个 regime。第一，数据与 stem 存在性论题一致：一个不受限
的树搜索——可以返回任何收缩树——\emph{从未}超越 stem 线一个语义换算单位以上，这正是"晶格
网络的最优收缩可取 stem 结构"所预言的；那两个 $-1$ 点是换算残差，不是反驳。第二，在八个实
例上（$8\le N\le 28$，外加 $N=48$）找到的宽度与认证尺度数值重合——在两种语义许可的粒度
（$\pm1$）上一致，且
据我们所知，这是启发式收缩顺序第一次获得对其尺度的绝对确认，而非与其他启发式的比较。第三，
$N\in\{32,36,44,53\}$ 处固定预算的搜索（$128$ 次重复——刻意保守；生产设置高出几个量级）漂
到尺度\emph{之上} $+1$ 到 $+2$——在 $N=36$ 处返回的顺序与认证线宣告可达的水平差了张量尺寸
的 $4$ 倍，而只做相对评判的方法论永远不会标记它。漂移度量的是\emph{该预算下}的搜索质量，不
是 hyper-optimizer 的极限行为——而标尺让后续问题成为可量化的，于是我们量了：把漂移点在
$512$ 与 $4096$ 次重复下重跑，每个找到的宽度都回到认证线的换算容差之内
（$N\in\{32,44\}$：$512$ 下间隙 $0$；$N=36$：$-1$；$N=53$：$512$ 下 $+1$、$4096$ 下恰为
$54=N+1$），且没有任何点跌破该容差的下沿。随预算增长，找到的宽度收拢进认证尺度周围的换算
容差之内，且从未跌破其下沿。

\paragraph{从校准 grid 到 RQC 架构。}
grid 族负责校准标尺；stem 论题真正要生活的地方是随机量子电路的\emph{架构}，所以我们把同一
套流程跑在 RQC 族的门级时空网络上（单个封闭振幅，每个两比特门一个 4 阶张量）：深 regime 的
1D brickwork 链（$L=2n$ 层）、固定 $n=24$ 的深度过渡扫描、以及 Sycamore 式 ABCD 耦合花样的
2D 芯片~\cite{arute2019quantum}。收缩宽度只取决于网络图，所以这些正是随机电路的图；对每个实
例我们还\emph{展示}两条真实的 stem 序——时间扫描与空间扫描——并精确计算其宽度，使对比成为
"自由树搜索 vs.\ 理论开出的 stem"（表~\ref{tab:rqc}）。

\begin{table}[t]
\caption{RQC 架构：找到的最优树宽度 vs.\ 两条展示 stem 扫描中较优者（每实例 $128$ 次重复）。
与图~\ref{fig:gapplot} 的 grid 族不同，此处的 "stem" 列是\emph{展示的上界}而非认证下界：
芯片晶格的宽度下界尚未形式化（\S\ref{sec:limits}），故 2D 各行比较的是搜索质量与一条具体
stem 的差距，而非与定理的差距。stem 论题是一个\emph{深 regime} 命题：其定义域是纠缠光锥覆盖全部比特之处，即可路由 bond 数
——由调度计数定理 \texttt{routed\_card\_le} 与 \texttt{exists\_schedule\_routed\_card} 钉在
$\min(n,b\cdot d)$（\S\ref{sec:routing}）——饱和到 $n$ 之处；浅层各行位于
\emph{该定义域之外}，列出只为描出 regime 的起点。}
\label{tab:rqc}
\centering\small
\begin{tabular}{llrl}
\hline
族 & 实例 & stem & found $-$ stem \\
\hline
1D 链，深（$L{=}2n$） & $n \in \{8, 12, 16\}$ & $n$ & $0,\ 0,\ 0$ \\
                      & $n \in \{20, 24, 28\}$ & $n$ & $+2,\ +2,\ +2$ \\
1D 链，极深（$L{=}3n$） & $n \in \{16, 20, 24, 28\}$ & $n$ & $0,\ +2,\ +4,\ +4$ \\
1D 链，$n{=}24$，深度 $L$ & $L \in \{6, 10, \ldots, 22\}$（阈值之下） & $\sim L{+}1$ & 各 $-2$（域外） \\
                      & $L \in \{26, 30, 38, 46\}$（深） & $24$ & $0,\ +1,\ +2,\ +4$ \\
2D 芯片，ABCD & $16$--$49$ 比特，$8$--$20$ 周期（$11$ 实例） & $n$ & $+2$ 到 $+17$，随规模增长 \\
\hline
\end{tabular}
\end{table}

\textbf{全部二十五个深 regime 实例中，没有任何找到的树胜过所展示的 stem}
——stem 论题在真实 RQC 图上、恰在其声明的定义域上受检。这个定义域在物理上和形式上都有含
义：阈值之下纠缠光锥尚未把所有比特连成一体，割面由光锥宽度而非 $n$ 支配，此时下界与论题都不
在受检状态——那正是 Napp 的可模拟相~\cite{napp2022efficient}，而其边界由
\S\ref{sec:routing} 的光锥 walk 定理标出（\texttt{bond\_causal}：显式 walk 认证第 $i$ 对
跨割面比特在 $i<d$ 时可达；逆方向未形式化）。浅层的 $-2$ 项位于定义域之外，其次才是在那里衡量两条典范见证的松弛。接着，
$n{=}24$ 的深度扫描干净地描出进入深 regime 的过程：交叉点恰在空间扫描超过时间扫描处
（$L\approx 26\approx n$），即光锥首次跨越全链之处；而一旦饱和，stem 尺度便与深度无关——
把 $L{=}2n$ 加深到 $3n$，stem 仍停在 $n$，搜索却漂得更远。最后，在
Sycamore 类 2D 架构上，自由搜索返回的树比平凡的时间扫描 stem 高出 $+2$ 到 $+17$ 个单位，且
缺口\emph{随规模增长}：$7\times 7$ 芯片（$49$ 比特，Sycamore-53 的尺寸级别）上，返回顺序物
化的张量比理论免费奉上的 stem 大 $2^{17}\approx 10^{5}$ 倍。在 supremacy 类架构上，搜索者
未能\emph{找到} stem，且缺口恰恰在最要紧的实例上加剧。更大
的预算也救不回来：在 $512$ 与 $4096$ 次重复（$4\times$ 与 $32\times$）下，$6\times6$ 芯片
的缺口只缩到 $+10$、再到 $+6$，$7\times7$、$12$ 周期实例缩到 $+16$/$+10$——预算放大
$32\times$ 之后仍是 stem 张量尺寸的 $2^{6}$--$2^{10}$ 倍，收拢速度约为每 $8\times$ 预算
$4$ 个宽度单位。就预算而言，搜索只能以对数速度收拢到理论免费奉上的 stem——这为"认证的
stem 感知排序"提供了一条直接的实践论据。

%% ===================================================================
\section{几何路由：brickwork 与光锥}\label{sec:routing}

\S\ref{sec:rigidity} 的刚性证书消耗一个跨割面匹配，怀疑者理应追问：真实硬件是否真能供给所
需尺寸的匹配。开发用两条关于 brickwork 电路的组合定理作答，把物理挡在 trusted base 之外。
其一是计数封顶：深度 $d$、每层至多 $b$ 个跨割面门的调度，至多路由 $b\cdot d$ 个不同的 bond
——是证出来的，不是假设的（\texttt{routed\_card\_le}，一个 \texttt{card\_biUnion} 界）——
且 $\min(n, b\cdot d)$ 由显式调度达到。其二是因果可达：bond 不是抽象标签，而是必须处于彼此
深度 $d$ 光锥内的比特对。在链上，一条长度 $2i+1$ 的显式 walk 连接第 $i$ 个嵌套跨割面对，
故当 $i<d$ 时该对处于深度 $d$ 的光锥之内（\texttt{bond\_causal}；逆方向——$i\ge d$ 时不可
达，即耦合图距离下界——未形式化）；在二维芯片上，链 walk 逐行经盒积结构抬升，在方形芯片上给出尺寸
$\min(n/2, \sqrt{n}\,d)$ 的因果可实现匹配——又是 $\mathrm{sweepCut}$ 尺度。合起来，这两条
为物理路由一步\emph{划定范围}：深 brickwork 布局的耦合图中确实存在所需尺寸的跨割面配对，且
每对都在光锥距离之内。两条限定为其划界。其一，电路动力学模型未被形式化——"深度 $d$ 的
电路能纠缠其光锥内的一对比特"是建模对应，上述 walk 定理是关于耦合图距离的陈述。其二，
\S\ref{sec:endtoend} 的刚性见证并不消费这些定理：它们直接以块对齐匹配为数据。物理芯片具体
的 $(2{+}1)$ 维 ABCD 耦合调度是我们尚未形式化的进一步保真度。

%% ===================================================================
\section{相关工作}\label{sec:related}

\paragraph{TNCO 实践。}hyper-optimized 收缩树~\cite{gray2021hyper}、big-batch 收
缩~\cite{pan2022solving}、stem 序~\cite{huang2020classical}、multi-stem 调
度~\cite{swtnc2025}、复用感知搜索~\cite{kalachev2021recursive}定义了当前最前沿；电路图的
treewidth 估计可追溯到 QuickBB 风格的启发式~\cite{boixo2017simulation}。这条线上的一切质量
评判都是相对的；没有一个产出对最优的证书。唯一的例外恰恰印证了规律：穷举动态规划对小网络返
回保证最优的收缩序~\cite{pfeifer2014faster}，但该保证落在搜索实现上，且方法在实例开始要紧
之处就止步了。我们的工作与之互补：我们不搜索顺序；我们从下方界住搜索所返回顺序的代价。

\paragraph{宽度理论。}收缩--treewidth 对应是 Markov--Shi~\cite{markov2008simulating}；
grid bramble 归于 Seymour 与 Thomas~\cite{seymour1993graph}，多维 grid 的值由
Kozawa--Otachi--Yamazaki~\cite{kozawa2014treewidth} 给出，pathwidth--treewidth 关系由
Groenland 等人~\cite{groenland2023tight}量化。求解器一侧，PACE 挑战一系的精确 treewidth
实现如今会给其下界附上基于 minor 的紧凑证书~\cite{tamaki2022heuristic}；这些证书由未经验
证的代码检查——这正是我们的 bramble 验证所补上的信任缺口。我们在此的贡献不是新图论，而是
其 pathwidth 下界机器的首个认证形态，并且工程化到只需要区间 Helly。

\paragraph{对宽度范式的逃逸。}decision diagram 在结构化取值上胜过宽
度~\cite{dd2025treewidth}，浅层随机 2D 电路可被高效模拟~\cite{napp2022efficient}。两者都塑
造了我们的陈述：压缩下界显式地限定为 \emph{generic}，regime 假设生活在路由层。

\paragraph{形式化。}证明助理中的图论先例包括 Doczkal 与 Pous 的 Coq
库~\cite{doczkal2020graph}（minor、treewidth-2 刻画），以及 Dittmann 的 Isabelle/HOL 树分
解~\cite{dittmann2016tree}（证到树与完全图的 treewidth）；两者都未形式化 pathwidth、结构化
图的对偶证书、或任何与收缩的联系。张量一侧，TNLean 形式化了矩阵乘积态的基
本定理~\cite{tnlean2026}，同样不涉及收缩复杂度。据我们所知，这是任何证明助理中首个
pathwidth 下界理论与首个 TNCO 开发。精神上我们的证书追随 verified SAT/LP checking——
cake\_lpr 用 verified 检查器验证不受信的 UNSAT 证明~\cite{tan2021cakelpr}——搜索不受
信、验证是 verified 的，而 bramble 对宽度扮演的正是对偶解对 LP 下界的角色。

%% ===================================================================
\section{局限与未来工作}\label{sec:limits}

证书适用范围，按用户可以依赖的声明来陈述。(1)~无条件下界量化在 stem（sweep）序上；推广到任
意收缩树是 $\pw=\Theta(\tw)$ 晶格输入，而开发中不存在树收缩序的类型——该条件性树陈述只生
活在 prose 中，其具名 \texttt{Prop} 形式在我们的模型上真空成立（序空间全是 stem）。形式化
子树 Helly 与树分解将使之成为定理，是自然的下一步。今天的证书读作：没有任何 sweep 型顺序低
于 $54$；条件性地、且在形式化之外，没有任何收缩树低于 $\Theta(54)$。(2)~实例化的晶格是深 regime 线链时空 grid；$(2{+}1)$ 维芯片晶格的
bramble~\cite{kozawa2014treewidth} 是后续工作，而对一般不规则网络框架已经适用——任何
bramble，无论如何找到，都产出认证下界——因此为实用网络整理 bramble 库是工具问题而非理论缺
口。(3)~\texttt{orderCost} 数的是物化的边界张量分量；常数因子的 FLOP 模型与切片权衡位于该
下界之上（切片以内存换时间，不减少总功）。(4)~压缩下界按设计是 generic 的：结构化实例
（stabilizer、真正低秩）可以也应当击穿它，而离散门集上的具体电路不被真子簇陈述覆盖；它也是
一张独立证书、陈述于主定理合取之外——没有任何定理把它与 \texttt{orderCost} 下界复合，且基
准见证族在门参数上是常量（\S\ref{sec:endtoend}）。
(5)~零测度推论（$\Pr[\text{坍缩}]=0$）已陈述、尚未形式化。(6)~stem 收缩序与路径分解的等同
（\S\ref{sec:background}）、\texttt{orderCost} 与收缩功的等同、以及完整 $N\times N$ grid
图与深 brickwork 链更稀疏的门级网络图的等同，都是以 prose 论证的建模对应；形式化的序空间
\emph{就是}该 grid 图的路径分解。(7)~实证对比（\S\ref{sec:empirical}）报告固定搜索预算，并附 $128/512/4096$ 次重复的敏感
性扫描：grid 漂移随预算增长收拢到换算容差之内，2D 芯片缺口则持续存在（$4096$ 下仍
$+6$ 到 $+10$）——观测到的间隙度量所述预算下的搜索质量，非极限行为。(8)~路由脚手架只认证
光锥边界的一个方向：显式 walk 见证 $i<d$ 时的可达性（\texttt{bond\_causal}）；逆方向的距
离下界（$i\ge d$ 时不可达）未形式化。

%% ===================================================================
\section{结论}\label{sec:conclusion}

TNCO 搜索如今有了一把绝对标尺：一个机器核验的代价下界——bramble 对偶证书管宽度、显式代价
语义、旁侧另有一条秩一侧的框架层面 generic 无压缩定理（其基准实例化仅认证能力，
\S\ref{sec:endtoend}）——公理足迹恰为 \texttt{propext}、
\texttt{Classical.choice}、\texttt{Quot.sound}。在校准算例上标尺在宽度上是精确的：认证下界
在 $54$ 处与显式 sweep 相遇，一份 stem 类内的机器核验最优性证明。余下的工作扩展证书的适用
范围而不触动其信任基础：经子树 Helly 到树序、到芯片晶格、再到启发式至今盲飞的不规则网络的
bramble 库。

\appendix
\section{Artifact（对照英文版附录）}\label{app:artifact}

匿名 artifact 档案（链接经投稿系统提供）包含完整 Lean 开发与实证流水线；档案中的 Lean 命名
空间与模块前缀亦由打包脚本改为匿名名称。Lean 开发：18 个
模块（17 个证明模块加公理审计模块）、约 2{,}800 行，\texttt{lake build} 构建于钉定
v4.30.0 的 Lean~4 / Mathlib；无
\texttt{sorry}、无 \texttt{axiom} 声明；检查模块
（\texttt{CheckAxioms.lean}，由库根引入、随 \texttt{lake build} 构建）收集
本文引用的每条定理的公理依赖，若任一集合超出
\texttt{\{propext, Classical.choice, Quot.sound\}} 则\emph{构建失败}
（审计断言的是包含关系——安全方向；正文的"恰好"是对每条被引定理的实测值）。
论文—artifact 对照：\S\ref{sec:bramble} \texttt{Bramble} / \texttt{GridConn}；
\S\ref{sec:exact} \texttt{GridExact}；\S\ref{sec:model} \texttt{GridModel}；
\S\ref{sec:rigidity} \texttt{Rigidity} / \texttt{Matching}；\S\ref{sec:routing}
\texttt{Brickwork} / \texttt{Causal}；\S\ref{sec:endtoend} \texttt{Sycamore}。实证流水
线：cotengra 脚本、Slurm 批处理文件与 \S\ref{sec:empirical} 背后的原始 JSON 结果（含预算
敏感性运行），溯源记录于实验 README（集群名为匿名而隐去）。

\bibliographystyle{plain}
\bibliography{references}

\end{document}
